% Propietario conservado: /Users/ruben/Documents/ChatGPT/jueces y controles/REVISION_ARGUMENTAL_APERTURAS_CIERRES_20260915/II/manuscript_es/sections/08j_enlace_angular_witt.tex
% FORMALIZACION_REUNIDA: alineamiento dodecafasico y lector angular.
% CERTIFICADO_NUEVO: composicion exacta conservada en el suplemento local.
% La generacion previa del sello y la seleccion sectorial mantienen sus pruebas.
\section{Equivariant realization of the angular reader in Witt incidence}
\label{ii:sec:ii-enlace-angular-witt}

The angular chart and exceptional incidence preserve different information from the same APP--TRIT--TPK development. The former expresses rotation coordinates; the latter preserves relations between positions, supports and orientations. Comparing them requires a transport that preserves these relations and allows the complete coordinates to be recovered. We construct that transport here and prove its compatibility with oriented permutations of the regional triplet.

The antecedents of the composition are the dodecaphasic seal \(K\), its marked face \(B_K\), the sector of the three-dimensional isotypic projector \(P_3^{\mathrm{ico}}\), the compatible branches of the nonadic filtration and the sectorial system of Definition~\ref{ii:def:ii09-sistema-sectorial}. Each enters with its structure, not merely through its scalar value. The real vector spaces of this subsection are subsequent realizations of those already expressed sectors. Generating the seal, transporting its coordinates and evaluating the reader correspond to distinct stages of the causal composition.

\subsection{Regional triplet and orientation line}

Let \(V\) be the counting space of the three positive regions of
\(\pi_{\rm HMT}\), with ordered orthonormal basis
\[
 (v_1,v_2,v_3)=(933555,393555,339555).
\]
The notation identifies each basis vector with the region it represents.
In the prolongation window of index seven, the eleven admitted branches
are distributed into orbits of sizes \(3+3+1+1+3\) under simultaneous
permutation of the final three positions of the first two words.
The sector retaining continuation to horizon two consists of
\[
\begin{array}{c|c|c}
\text{branch}&\text{triple of words}&\text{positive region}\\ \hline
0&200112\mid212221\mid110110&339555\\
1&200121\mid212212\mid110110&393555\\
2&200211\mid212122\mid110110&933555.
\end{array}
\]
The central branch additionally extends to horizon three. These
horizons belong to the same filtration of compatible prolongations;
the successive support cardinalities are \(11\), \(3\), and \(1\).

The map from the tail of the first word to the positive region
substitutes \(1\mapsto3\) and \(2\mapsto9\) coordinatewise and adds
the regional tail \(555\). It therefore preserves the position of the
exceptional symbol and commutes with every permutation of the three positions. The
branch order is \((339555,393555,933555)\); the change from the order of
\(V\) is the permutation \(R=(1\,3)\).

The orientation line of \(V\) is \(o(V)=\det V\). The map
\begin{equation}
 \Psi:\Lambda^2V\longrightarrow\operatorname{Hom}(V,\det V),
 \qquad \Psi(u\wedge v)(w)=u\wedge v\wedge w
 \label{ii:eq:ii-witt-hodge-natural}
\end{equation}
is identified through the counting metric with an isometry
\(\Lambda^2V\to o(V)\otimes V\). For
\(\omega=v_1\wedge v_2\wedge v_3\), its values on the exterior basis
are
\[
 v_1\wedge v_2\mapsto\omega\otimes v_3,
 \qquad v_2\wedge v_3\mapsto\omega\otimes v_1,
 \qquad v_1\wedge v_3\mapsto-\omega\otimes v_2.
\]
In angular order \((12,23,13)\), joint transport of basis and
orientation has matrix
\begin{equation}
 D=\det(R)P_R
 \begin{pmatrix}0&1&0\\0&0&-1\\1&0&0\end{pmatrix}
 =\begin{pmatrix}-1&0&0\\0&0&1\\0&-1&0\end{pmatrix}.
 \label{ii:eq:ii-witt-orientation-D}
\end{equation}
Thus
\(D(\theta_{12},\theta_{23},\theta_{13})^{\mathsf T}
=(-\theta_{12},\theta_{13},-\theta_{23})^{\mathsf T}\).
Throughout this section, \(P_p e_j=e_{p(j)}\) defines the permutation-matrix
convention, and the asterisk denotes the adjoint for the declared
counting metrics.

\begin{proposition}[Naturality of the oriented identification]
\label{ii:prop:ii-witt-D-natural}
For every \(p\in S_3\),
\[
 D\Lambda^2P_p=\det(p)P_{RpR^{-1}}D,
 \qquad D^*D=I_3.
\]
\end{proposition}
\begin{proof}
The identity
\((P_pu)\wedge(P_pv)\wedge(P_pw)=\det(p)(u\wedge v\wedge w)\)
proves naturality of \eqref{ii:eq:ii-witt-hodge-natural}. Simultaneous
change of regional order introduces conjugation by \(R\).
Finally, the columns of \(D\) are signed coordinate vectors,
whence \(D^*D=I_3\).
\end{proof}

The orientation line explains the difference between the character
\((3,1,0)\) of the position triplet and the character \((3,-1,0)\)
of the oriented triplet, on the identity, transposition, and
order-three cycle classes. The three transpositions provide an explicit control:
omitting the determinant factor changes the sign of the preceding identity.
A global change of the oriented anchor is instead transported as a joint
coordinate change.

\subsection{Dodecaphase sector and variational selection of orientation}

The chart of the previously expressed seal is
\begin{equation}
 K=(234,543,140,729,659,824,621,58,914,794,146,601).
 \label{ii:eq:ii-witt-K-marked}
\end{equation}
To make its incidence realization explicit, write the ternary Paley
matrix in the gauge used:
\[
 A_W=\begin{pmatrix}
 0&1&1&1&1&1\\
 1&0&1&2&2&1\\
 1&1&0&1&2&2\\
 1&2&1&0&1&2\\
 1&2&2&1&0&1\\
 1&1&2&2&1&0
 \end{pmatrix},
 \qquad C(a)=(a,aA_W),\quad a\in\mathbb F_3^6.
\]
Support is taken on the twelve coordinates of \(C(a)\). The
marked intersection has the two compatible expressions
\begin{equation}
 B_K=\operatorname{supp}C(010211)\cap\operatorname{supp}C(201101)
     =\{p:K_p\ge729\}=\{4,6,9,10\}.
 \label{ii:eq:ii-witt-B-marked}
\end{equation}
The equality is evaluated by ternary multiplications of the displayed
matrix. In particular, the intersection is obtained as a set of
positions before its cardinality four is used.

To fix the representation of the twelve positions completely,
consider \(A_5\) acting by left multiplication on the
cosets \(r_pC_5\), where \(C_5=\langle(1\,2\,3\,4\,5)\rangle\).
The representatives, written as image lists on
\(\{1,\ldots,5\}\), are
\[
\begin{array}{c|c@{\qquad}c|c}
p&r_p&p&r_p\\ \hline
1&(4,3,5,2,1)&7&(5,4,3,2,1)\\
2&(4,2,1,3,5)&8&(1,3,4,2,5)\\
3&(1,2,4,5,3)&9&(4,2,3,5,1)\\
4&(2,5,3,4,1)&10&(3,2,5,4,1)\\
5&(4,3,1,5,2)&11&(4,5,2,3,1)\\
6&(5,1,2,3,4)&12&(5,3,2,4,1).
\end{array}
\]
Denote this permutation representation by \(\rho\), and put
\(V_{11}=\mathbf1^\perp\subset\mathbb R^{12}\). The icosahedral
three-dimensional isotypic projector is determined by
\begin{equation}
 P_3^{\mathrm{ico}}=\frac3{60}\sum_{a\in A_5}\chi_3(a^{-1})\rho(a),
 \label{ii:eq:ii-witt-projector-character}
\end{equation}
where the values of \(\chi_3\) on classes of types
\(1,2^2,3,5_a,5_b\) are
\[
 \left(3,-1,0,\frac{1+\sqrt5}{2},\frac{1-\sqrt5}{2}\right).
\]
Class \(5_a\) contains \((1\,2\,3\,4\,5)\), and class
\(5_b\) contains its square. This convention distinguishes the two
three-dimensional summands. The superscript \(\mathrm{ico}\) identifies
this component of the \(A_5\) representation; the cyclic projector
\(P_3\) entering~\eqref{ii:eq:det} acts on twelve-position transport
with its own definition. Finite evaluation of
\eqref{ii:eq:ii-witt-projector-character} satisfies
\(P_3^{\mathrm{ico}}=(P_3^{\mathrm{ico}})^*=(P_3^{\mathrm{ico}})^2\), \(P_3^{\mathrm{ico}}\mathbf1=0\), and
\(\operatorname{Tr}P_3^{\mathrm{ico}}=3\).

The two marked vectors of the sector are
\[
 b=P_3^{\mathrm{ico}}\left(\mathbf1_{B_K}-\frac13\mathbf1\right),
 \qquad u=P_3^{\mathrm{ico}}(K-\overline K\mathbf1),
 \qquad \overline K=\frac1{12}\sum_{p=1}^{12}K_p,
 \qquad \|b\|^2=1.
\]
For each of the ten order-three subgroups \(C_3\subset A_5\),
let \(H=N_{A_5}(C_3)\cong S_3\). The character
\(\operatorname{sgn}_H\) takes value \(+1\) on \(C_3\) and
\(-1\) on \(H\setminus C_3\). It is the character of the quotient
\(H/C_3\), whereas the ambient sign of elements of \(A_5\)
is always positive. Define
\begin{equation}
 e_{{\rm sgn},H}=\frac16\sum_{a\in H}\operatorname{sgn}_H(a)\rho(a),
 \qquad \mathcal E_B(H)=\|e_{{\rm sgn},H}b\|^2.
 \label{ii:eq:ii-witt-energy}
\end{equation}

\begin{proposition}[Normalizer and orientation determined by the seal]
\label{ii:prop:ii-witt-variational-selection}
The maximization rule of \eqref{ii:eq:ii-witt-energy} selects a unique
normalizer \(H_*\). Its energy and the second distinct energy are
\[
 \frac23+\frac{2\sqrt5}{15},\qquad
 \frac13+\frac{2\sqrt5}{15},
\]
respectively; the gap is \(1/3\). Projection of \(u\)
selects the same normalizer. Within \(H_*\), maximizing
\(s(a)=\langle b,\rho(a)u\rangle\), separately among order-three
elements and involutions, determines
\[
 c_*=(1,2,4,5,3),\qquad r_*=(2,1,5,4,3)
\]
in image notation on five points. The identity
\(r_*c_*r_*=c_*^{-1}\) holds.
\end{proposition}
\begin{proof}
Evaluation of the preceding finite sums in \(\mathbb Q(\sqrt5)\)
gives the following distribution of the ten energies:
\[
\begin{array}{c|c@{\qquad}c|c}
\text{energy}&\text{multiplicity}&\text{energy}&\text{multiplicity}\\ \hline
2/3+2\sqrt5/15&1&2/3-2\sqrt5/15&1\\
1/3+2\sqrt5/15&2&1/3-2\sqrt5/15&2\\
1/6+\sqrt5/30&2&1/6-\sqrt5/30&2.
\end{array}
\]
The greatest belongs to the subgroup generated by \(c_*\). The next is
obtained by subtracting \(1/3\). The same evaluation with \(u\) has a unique
maximum at that subgroup, of value
\(2526869/12+2813609\sqrt5/30\).
The scores of the two cycle orientations are
\[
 s(c_*)=\frac{1341}{4}+\frac{2409\sqrt5}{20},\qquad
 s(c_*^{-1})=\frac{357}{2}+\frac{1853\sqrt5}{10}.
\]
Their difference is \(627/4-1297\sqrt5/20>0\), as is checked
by comparing positive squares. The scores of the three
involutions are
\[
\begin{array}{c|c}
\text{list of images}&s(a)\\ \hline
(2,1,3,5,4)&-1341/4-2409\sqrt5/20\\
(2,1,4,3,5)&-1659/4-2667\sqrt5/20\\
(2,1,5,4,3)&-357/2-1853\sqrt5/10.
\end{array}
\]
The unique maximum corresponds to \(r_*\). Composition of the two image
lists verifies the dihedral relation. All these operations use
\(K\), \(B_K\), \(P_3^{\mathrm{ico}}\), and the finite actions already specified.
\end{proof}

Write \(\vartheta:S_3\to H_*\) for the isomorphism determined
by \(\vartheta((1\,2\,3))=c_*\) and
\(\vartheta((1\,3))=r_*\). The position reflection fixes the central
branch reaching horizon three. This orientation completes
regional identification before any angle evaluation.

\subsection{Reynolds intertwiner and polar normalization}

In the branch basis, let \(\rho_o(g)=\det(g)P_g\) and
\(w=(0,1,0)^{\mathsf T}\). The Reynolds operator defined by summation
over the group is
\begin{equation}
 A_R=\sum_{g\in S_3}\rho_o(g)\,|w\rangle\langle b|\,
                       \rho(\vartheta(g))^{-1}.
 \label{ii:eq:ii-witt-reynolds}
\end{equation}
The symbol \(A_R\) distinguishes this operator from the angle
\(A_{\rm deg}\). Each summand has unit scale; metric
normalization is performed through its polar factor.

\begin{proposition}[Polar isometry of the oriented sector]
\label{ii:prop:ii-witt-polar}
The following hold
\[
 A_R=A_RP_3^{\mathrm{ico}},\qquad
 \rho_o(g)A_R=A_R\rho(\vartheta(g)),\qquad
 G=A_RA_R^*>0,\qquad A_R^*G^{-1}A_R=P_3^{\mathrm{ico}}.
\]
With \(E=\mathbf1\mathbf1^{\mathsf T}/3\),
\begin{equation}
 G=\lambda_sE+\lambda_t(I_3-E),\qquad
 \lambda_s=8+\frac{8\sqrt5}{5},\quad
 \lambda_t=\frac12-\frac{\sqrt5}{5}>0.
 \label{ii:eq:ii-witt-gram-spectrum}
\end{equation}
Consequently,
\[
 U_{GK}=G^{-1/2}A_R:P_3^{\mathrm{ico}}V_{11}\longrightarrow o(V)\otimes V
\]
is a surjective isometry:
\(U_{GK}^*U_{GK}=P_3^{\mathrm{ico}}\) and \(U_{GK}U_{GK}^*=I_3\).
\end{proposition}
\begin{proof}
The projector \(P_3^{\mathrm{ico}}\) commutes with \(\rho\), and \(b=P_3^{\mathrm{ico}}b\);
hence \(A_R=A_RP_3^{\mathrm{ico}}\). Reindexing the Reynolds sum proves
equivariance. Multiplication of the finite sum gives diagonal entries
\(3+2\sqrt5/5\) and off-diagonal entries
\(5/2+3\sqrt5/5\) in \(G\). The projections \(E\) and \(I_3-E\)
give \eqref{ii:eq:ii-witt-gram-spectrum}. Both eigenvalues are
positive; for the smaller, \(25>20\) suffices. Thus \(A_R\) has
rank three. Its row space is contained in the rank-three sector
of \(P_3^{\mathrm{ico}}\), and equals it. The orthogonal projector onto that row
space is \(A_R^*(A_RA_R^*)^{-1}A_R\), giving the final identity.
The two polar-factor equations follow by substitution.
\end{proof}

The positive root preserves equivariance because both spectral
projectors commute with the oriented representation. Uniqueness of the
polar factor pertains to the Reynolds operator and orientation
determined by the preceding rules.

\subsection{Point--hexad incidence and joint gauge change}

Let \(\mathcal W\) be the design \(S(5,6,12)\) of 132 hexads obtained
from the weight-six supports of the ternary code \(C(a)\). The incidence
matrix is
\[
 I_{H,p}=\mathbf1_{p\in H},\qquad
 (Iz)_H=\sum_{p\in H}z_p.
\]
The blocks are reconstructed by enumerating \(a\in\mathbb F_3^6\),
retaining weight six, and collecting equal supports. The finite
verification that each five-point subset belongs to a unique
hexad fixes the design property used below.

\begin{lemma}[Metric and spectral incidence identities]
\label{ii:lem:ii-witt-incidence-identities}
Let \(J_{12}\) and \(J_{132,12}\) be the all-ones matrices indicated by
their dimensions, and let
\[
 (G_{\mathcal W})_{H,H'}=\binom{|H\cap H'|}{4}.
\]
Then
\begin{equation}
 I^*I=36I_{12}+30J_{12},\qquad
 G_{\mathcal W}I=30I+15J_{132,12}.
 \label{ii:eq:ii-witt-incidence-identities}
\end{equation}
In particular, \(U_W=I/6\) is an isometry on
\(V_{11}=\mathbf1^\perp\), with image contained in
\(E_{30}=\ker(G_{\mathcal W}-30I_{132})\).
\end{lemma}
\begin{proof}
The design property gives 66 hexads per point and 30 per pair of points:
\[
 \frac{\binom{11}{4}}{\binom54}=66,
 \qquad \frac{\binom{10}{3}}{\binom43}=30.
\]
These are the diagonal and off-diagonal entries of \(I^*I\).
For the second identity, count pairs \((T,H')\), where
\(T\subset H\cap H'\), \(|T|=4\), and \(p\in H'\).
If \(p\in H\), ten tetrads \(T\subset H\) contain \(p\),
and each belongs to four hexads; the other five tetrads,
adjoined to \(p\), determine one hexad each. The result is
\(10\cdot4+5=45\). If \(p\notin H\), each of the fifteen
tetrads determines a hexad together with \(p\), and the result is 15.
This yields \(30I_{H,p}+15\). The all-ones matrices vanish on
\(V_{11}\), proving the isometry and spectral membership.
\end{proof}

The gauge change of the twelve positions is
\begin{equation}
 \sigma=(1,5,2,6,7,4,3,10,11,9,8,12),
 \label{ii:eq:ii-witt-sigma}
\end{equation}
written as an image list with indices starting at one. If
\(\rho_{\rm perm}\) is the point action underlying \(\rho\), the
six permutations
\[
 \tau_g=\sigma\rho_{\rm perm}(\vartheta(g))\sigma^{-1},\qquad g\in S_3,
\]
preserve \(\mathcal W\). Their two generators have image lists
\[
\begin{aligned}
 \tau_{(123)}&=(3,6,5,2,1,4,9,12,11,8,7,10),\\
 \tau_{(13)}&=(2,1,4,3,6,5,8,7,10,9,12,11).
\end{aligned}
\]
The assertion is checked on the finite design: these two generating
permutations are applied to the weight-six supports of
\(C(a)\), and their images are verified to run through the
same 132 supports. Products of these generators give the other
four permutations. The induced action \(Q_{\tau_g}\) on hexads
satisfies
\begin{equation}
 I P_{\tau_g}=Q_{\tau_g}I,
 \qquad
 P_\sigma\rho(\vartheta(g))=P_{\tau_g}P_\sigma.
 \label{ii:eq:ii-witt-recalibration}
\end{equation}
The first identity expresses exactly the equivalence
\(p\in H\Longleftrightarrow\tau_g(p)\in\tau_g(H)\).

This change jointly transports positions, code, incidence,
origin, chamber, and orientation of the seal. In the initial gauge, five
elements of this order-six subgroup do not preserve the fixed design.
Conjugation \eqref{ii:eq:ii-witt-sigma} provides the compatible gauge
and maintains that control over the initial labeling. The procedure
exhibits a gauge change; its uniqueness does not enter the proof.

\subsection{Isometric composition and coordinate recovery}

\begin{theorem}[Dodecaphase--Witt angular transport]
\label{ii:thm:ii-witt-angular-transport}
The composition
\begin{equation}
 \boxed{F=\frac16 I P_\sigma A_R^*G^{-1/2}D:
        \mathbb R^3_{\rm ang}\longrightarrow E_{30}}
 \label{ii:eq:ii-witt-F}
\end{equation}
is an isometry. Its coordinate action and equivariance are
\begin{align}
 (F\boldsymbol\theta)_H
 &=\frac16\sum_{p\in H}
       (P_\sigma A_R^*G^{-1/2}D\boldsymbol\theta)_p,
       \label{ii:eq:ii-witt-F-coordinates}\\
 F X_g&=Q_{\tau_g}F,
 \qquad X_g=\Lambda^2P_{R^{-1}gR}.
       \label{ii:eq:ii-witt-F-equivariance}
\end{align}
\end{theorem}
\begin{proof}
The columns of \(A_R^*\) belong to \(P_3^{\mathrm{ico}}V_{11}\), and
\(P_\sigma\) preserves \(V_{11}\).
Lemma~\ref{ii:lem:ii-witt-incidence-identities} therefore proves
membership of the image of \(F\) in \(E_{30}\). Moreover,
\begin{align*}
 F^*F
 &=D^*G^{-1/2}A_RP_\sigma^*
       \frac{I^*I}{36}P_\sigma A_R^*G^{-1/2}D\\
 &=D^*G^{-1/2}A_RA_R^*G^{-1/2}D
  =D^*D=I_3.
\end{align*}
The coordinate formula expands multiplication by incidence.
For equivariance, compose in order the naturality of \(D\),
equivariance of \(A_R^*\), commutation of \(G^{-1/2}\) with
\(\rho_o\), and the two identities of
\eqref{ii:eq:ii-witt-recalibration}. This produces
\eqref{ii:eq:ii-witt-F-equivariance} for each \(g\in S_3\).
\end{proof}

\begin{proposition}[Covariance of the complete gauge]
\label{ii:prop:ii-witt-gauge-covariance}
Let \(S,T,V_o\) be orthogonal coordinate changes in the dodecaphase
space, Witt point space, and oriented sector,
respectively. Let \(Q_T\) be the corresponding coordinate change
in the incidence space; for a design relabeling, it is the
induced hexad permutation. If one jointly transports
\[
\begin{gathered}
 A_R'=V_oA_RS^*,\quad G'=V_oGV_o^*,\quad
 P_\sigma'=TP_\sigma S^*,\\
 D'=V_oD,\qquad U_W'=Q_TU_WT^*,
\end{gathered}
\]
the composition satisfies \(F'=Q_TF\).
\end{proposition}
\begin{proof}
Uniqueness of the positive root gives
\((G')^{-1/2}=V_oG^{-1/2}V_o^*\). Substituting the five
transformations into
\(F'=U_W'P_\sigma'(A_R')^*(G')^{-1/2}D'\), the factors
\(T^*T\), \(S^*S\), and \(V_o^*V_o\) cancel in pairs, leaving
\(Q_TU_WP_\sigma A_R^*G^{-1/2}D=Q_TF\).
\end{proof}

The preceding covariance describes a coordinate change of all
marked data. Its application to an active cardinal transformation requires
the corresponding identity between actions on the state; the two
notions thus preserve their own domains.

\subsection{Faithful realization of the angular sectorial system}

We use the chart \(x_\theta\) of \eqref{ii:eq:ii-angular-chart},
\[
 x_\theta=\begin{pmatrix}A_{\rm deg}\\C^*_{\rm deg}/6\\
                              (180/\pi_{\rm HMT})\Delta_4\end{pmatrix}.
\]
The component \(h_\theta=C^*_{\rm deg}/6\) distributes the odd
component of the record among six antipodal pairs.
Definition~\ref{ii:def:ii09-sistema-sectorial} fixes the three rules
\[
 R_{12}=(2,-p,bq),\quad R_{23}=(0,q,-p^2/2),\quad
 R_{13}=(1,-pb,-(o_8-h_6)/t)
\]
for \((b,p,q,h_6,o_8,t)=(4,5,7,6,8,3)\). The face of cardinality
four, pentadic order, pivot coordinate, hexadic--octadic
difference, and tritic alphabet retain distinct types.
The bidirectional orbital normalization of
\ref{ii:prop:ii09-semicuadrado-orbital} preserves the term \(p^2/2\),
including the contribution of fixed pairs.

Its evaluation, already proved in
Theorem~\ref{ii:thm:ii09-matriz-sectorial}, gives
\[
 \mathsf M_{\rm CKM}=
 \begin{pmatrix}2&-5&28\\0&7&-25/2\\1&-20&-2/3\end{pmatrix},
 \qquad\det\mathsf M_{\rm CKM}=-\frac{3857}{6}.
\]
The transport \(F\) acts on the coordinates expressed by these rules; the preceding choice of functionals retains its place in the sectorial construction.

\begin{corollary}[Angular reader in the incidence module]
\label{ii:cor:ii-witt-composed-reader}
The map
\begin{equation}
 \mathcal C=F\mathsf M_{\rm CKM}
 \label{ii:eq:ii-witt-composed-reader}
\end{equation}
satisfies
\[
 \mathcal C^*\mathcal C=\mathsf M_{\rm CKM}^*\mathsf M_{\rm CKM},
 \qquad
 \mathsf M_{\rm CKM}^{-1}F^*\mathcal C=I_3.
\]
It is an isometry from the space of \(x_\theta\), equipped with metric
\(\mathsf M_{\rm CKM}^*\mathsf M_{\rm CKM}\), onto its image.
\end{corollary}
\begin{proof}
Substitute \(F^*F=I_3\) into both products. The nonzero determinant
of \(\mathsf M_{\rm CKM}\) permits use of its inverse,
displayed in \eqref{ii:eq:ii-ckm-inverse}. Positivity of the metric
follows from
\(x^*\mathsf M_{\rm CKM}^*\mathsf M_{\rm CKM}x
=\|\mathsf M_{\rm CKM}x\|^2>0\) for \(x\ne0\).
\end{proof}

The isometry \(F\) preserves the counting metric of
\(\boldsymbol\theta=\mathsf M_{\rm CKM}x_\theta\); its transport
to coordinates \(x_\theta\) gives the metric of the corollary. This
distinction avoids identifying a basis of angular generators with an
orthonormal basis of their parameters.

\subsection{Exact certification and finite unitary composition}

The positive roots of the polar factor can be checked through their
projectors while preserving exact arithmetic. Define
\begin{equation}
 B=\frac16IP_\sigma A_R^*D,\qquad
 H=D^*GD=B^*B,\qquad L=H^{-1}B^*.
 \label{ii:eq:ii-witt-radical-free}
\end{equation}
The identities
\[
 LB=I_3,\qquad L(B\mathsf M_{\rm CKM})=\mathsf M_{\rm CKM},\qquad
 \mathsf M_{\rm CKM}^{-1}L(B\mathsf M_{\rm CKM})=I_3
\]
are evaluated in \(\mathbb Q(\sqrt5)\). If \(E_D=D^*ED\),
\[
 H^{-1/2}=\lambda_s^{-1/2}E_D+
          \lambda_t^{-1/2}(I_3-E_D),\qquad F=BH^{-1/2}.
\]
The products \(E_D^2=E_D\), \(E_D(I_3-E_D)=0\), and the
decomposition of \(H\) prove this formula without approximating its
radicals. The 132 rows of \(B\) are determined by
\eqref{ii:eq:ii-witt-F-coordinates} and by the hexad construction
rule. The executable certificate of this delivery preserves those rows,
those of \(B\mathsf M_{\rm CKM}\), the left inverse \(L\),
the projectors, and the labels of each hexad. Its 906 exact
checks verify this composition and its orientation,
incidence, and recovery identities.

Conversion to radians is performed before the trigonometric realization:
\begin{equation}
 \boldsymbol\vartheta=
 \mathsf M_{\rm CKM}
 \begin{pmatrix}(\pi_{\rm HMT}/180)A_{\rm deg}\\
 (\pi_{\rm HMT}/180)C^*_{\rm deg}/6\\\Delta_4\end{pmatrix},
 \qquad
 \delta=\frac{\pi_{\rm HMT}}{180}\,9A_{\rm deg}.
 \label{ii:eq:ii-witt-angular-units}
\end{equation}
The finite realization preserves the order
\(V_{\rm CKM}=R_{23}U_{13}(\delta)R_{12}\) of
\eqref{ii:eq:ii-ckm-factorization}. To specify the additional information
of the complex phase, let \(E_{ab}\) be the elementary matrix and define
\[
 J_{13}=E_{13}-E_{31},\qquad S_{13}=E_{13}+E_{31},\qquad
 X_{13}(\delta)=\cos\delta\,J_{13}-i\sin\delta\,S_{13}.
\]
Then \(X_{13}(\delta)^*=-X_{13}(\delta)\) and
\(U_{13}(\delta)=\exp(\vartheta_{13}X_{13}(\delta))\).
Indeed, \(X_{13}(\delta)^2=-(E_{11}+E_{33})\), and its exponential
reproduces the unitary \(13\) block of
\eqref{ii:eq:ii-ckm-factorization}. Exterior transport preserves the
oriented generators; the symmetric imaginary component preserves
phase information. The ordered product of the three finite
factors maintains these two contributions and their order of composition.

\subsection{Structural result and scope of the composition}

The construction provides an equivariant isometric transport from the three angular coordinates to the Witt incidence, together with an explicit reconstruction of \(A_{\rm deg}\), \(C^*_{\rm deg}/6\) and \(\Delta_{\rm deg}\). It preserves the exterior orientation, the dodecaphasic sector, the marked face of the seal and the joint gauge change. In particular, the incidence realization expresses the same angular content on another support and allows its three coordinates to be recovered.

Prior generation of \(K\), selection of the sectorial
functionals, and physical identification of the amplitudes retain their
own proofs. The recovery identities proved here
establish faithfulness of the composition; its subsequent use keeps
those antecedents explicit. The complementary development of quartic
memory planes remains another construction of the state and does not
interpose itself as an additional requirement for the alignment realized in
\eqref{ii:eq:ii-witt-F}.
